\section{Related Work}
\label{sec:related}

\noindent \textbf{Specification-guided protocol analysis.}
Recent protocol-analysis systems increasingly use standard documents to guide
implementation checking. RIBDetector, RFCAudit, and ProtocolGuard locate
implementation inconsistencies or protocol non-compliance through RFC-rule
extraction, semantic code indexing, and LLM-guided program slicing with dynamic
validation~\cite{chen2022ribdetector,zheng2025rfcaudit,song2026protocolguard}.
ParCleanse, iPanda, and AUTOSPEC use LLMs, retrieval, execution analysis, or
implementation feedback to generate protocol message specifications,
quasi-oracles, formal specifications, or conformance
tests~\cite{zheng2025parcleanse,sun2025ipanda,liu2025autospec}. These systems
also use protocol standards as guidance, but their outputs are usually
rule-matching results, parser oracles, code snippets, or dynamic
counterexamples. In contrast, \tool{} models each check target as a field-level
conformance claim that connects the standard source, field semantics,
applicability condition, and implementation evidence. When the code context,
runtime path, or validation evidence is insufficient, the system preserves the
\textsc{Insufficient} state and retrieves additional evidence in the next
round, separating evidence-supported deviations from insufficiently supported
hypotheses. This combination of field-level modeling and non-binary judgment
allows \tool{} to characterize the correspondence between standard requirements
and implementation behavior at a finer granularity, especially for subtle
deviations in field values, state-dependent conditions, compound constraints,
and error handling. Our evaluation shows that \tool{} uncovers several
implementation deviations that, to the best of our knowledge, had not been
reported by prior standard-to-code checking work.

% The experimental results show that \tool{} can discover multiple implementation
% deviations that were not reported by prior standard-to-code checking work.

\noindent \textbf{Protocol testing.}
Traditional conformance suites and protocol fuzzers provide strong executable
evidence, but their oracles are usually hand-written, inferred from observed
behavior, or tied to protocol-specific testing frameworks. TLS-Attacker and
TLS-Anvil generate targeted TLS
tests~\cite{somorovsky2016tlsattacker,maehren2022tlsanvil}; protocol state
fuzzing, automata-based analysis, and monitor-based testing expose state-machine
deviations or conformance issues in TLS/DTLS, QUIC, and related protocol
implementations~\cite{de2015protocol,fiterau2020dtls,fiterau2023automata,beurdouche2015messy,asadian2024monitor,piraux2018quicevolution}.
Certificate-validation studies further use adversarial inputs, differential
testing, symbolic execution, and RFC-guided inputs to test SSL/TLS
implementations~\cite{brubaker2014frankencert,chen2015guided,chen2018rfc,chen2017hvlearn,chau2017symcerts},
while general protocol fuzzers such as AFLNet and StateAFL improve exploration
of stateful network
protocols~\cite{pham2020aflnet,natella2022stateafl,natella2021profuzzbench,ba2022stateful}.
These methods are effective at triggering implementation behavior and finding
defects, especially when test targets, input-generation strategies, or oracles
are already clear. In contrast, \tool{} focuses on identifying field-level
check targets from natural-language standards and linking them to
implementation code and runtime behavior, providing traceable standard
evidence for testing and human review.


\noindent \textbf{Formal verification.}
Prior work has explored how formal methods can be used to check conformance
between protocol specifications and implementations. Chudnov et al. use Cryptol
and SAW to prove that key components of Amazon s2n-TLS conform to high-level
specifications, and integrate the verification process into development
processes~\cite{s2ntls,cryptol,saw,chudnov2018s2n}. E-Verify constructs
Cryptol models of the handshake state machine from both the RFC and the OpenSSL
implementation, and compares state-transition sequences to check equivalence
between the standard and the
implementation~\cite{guan2025openssltls13}. SYMNV uses symbolic execution to
generate test inputs and checks whether network daemons violate specification
rules~\cite{song2011symnv}; EverParse and follow-up work generate verified
parsers and serializers from data-format
specifications~\cite{ramananandro2019everparse,swamy2022everparse3d}, while
recent work further explores end-to-end verification from high-level protocol
models to production code~\cite{pereira2024protocolstocode,arquint2022modular}.
Unlike formal verification methods that aim for full correctness proofs,
\tool{} avoids the high cost of manually constructing formal models and can
directly extract field-level requirements from natural-language standards and
locate implementation deviations.

% \noindent \textbf{Specification-to-code traceability and mining.}
% Requirements traceability research uses information retrieval, LSI,
% candidate-link generation, and deep semantic representations to connect
% natural-language requirements with source-code artifacts~\cite{antoniol2002recovering,marcus2003recovering,hayes2006advancing,clelandhuang2014software,guo2017semantically}.
% Specification mining infers likely properties from program executions or
% recurring code patterns~\cite{ernst2001dynamically,engler2001bugs,engler2001checking}.
% FASTRIC makes implicit FSM-like structures in LLM interactions explicit and
% checks execution traces against those structures~\cite{jin2025fastric}. 
% These lines of work primarily focus on link recovery, property inference, or
% execution-trace checking. In contrast, \tool{} treats standard--code linking as
% an intermediate step toward conformance judgments that include evidence and
% uncertainty notes, reducing the risk that incorrect links directly become false
% positives and improving the auditability of the results.
